\chapter*{CAPÍTULO III: ESPECIFICACIONES TÉCNICAS DE MATERIALES Y EQUIPOS}
\addcontentsline{toc}{chapter}{CAPÍTULO III: ESPECIFICACIONES TÉCNICAS DE MATERIALES Y EQUIPOS}

\section{ALCANCE Y NORMAS APLICABLES}

Las presentes especificaciones técnicas definen los requisitos de diseño, fabricación, pruebas en fábrica, transporte y recepción que deben cumplir los materiales y equipos electromecánicos destinados a la construcción de las redes de distribución primaria en 10.0 kV, subestaciones aéreas biposte (SAB) y redes de distribución secundaria en 380/220 V con alumbrado público en el casco urbano de Limbani.

Todos los suministros deberán cumplir rigurosamente con las siguientes normas técnicas:
\begin{itemize}
  \item \textbf{CNE - Suministro 2011:} Código Nacional de Electricidad - Suministro (R.M. N.° 366-2011-MEM/DM).
  \item \textbf{Normas DGE:} Resoluciones Directorales R.D. 018-2003-EM/DGE (Líneas Primarias), R.D. 016-2003-EM/DGE (Redes Secundarias) y R.D. 026-2003-EM/DGE (Especificaciones de Montaje).
  \item \textbf{NTP:} Normas Técnicas Peruanas del INACAL.
  \item \textbf{IEC / ANSI / ASTM:} Normas internacionales para equipos de transformación, aislamiento y protección contra sobretensiones.
\end{itemize}

\section{CONDUCTORES ELÉCTRICOS}

\subsection{Conductor de Media Tensión: Aleación de Aluminio AAAC 35 mm$^2$}
Para la línea primaria trifásica de 10.0 kV se empleará conductor desnudo de aleación de aluminio (Almelec / 6201-T81) conforme a la Norma Técnica Peruana NTP 370.258 y la norma IEC 61089:
\begin{itemize}
  \item \textbf{Sección nominal:} $35.0\text{ mm}^2$.
  \item \textbf{Número de alambres:} 7 alambres de diámetro uniforme ($2.52\text{ mm}$).
  \item \textbf{Diámetro exterior nominal:} $7.56\text{ mm}$.
  \item \textbf{Masa nominal aproximada:} $96.0\text{ kg/km}$ ($0.096\text{ daN/m}$).
  \item \textbf{Carga de rotura mínima certificada:} $10.50\text{ kN}$ ($1,070\text{ daN}$).
  \item \textbf{Resistencia eléctrica máxima a $20^\circ\text{C}$ (CC):} $0.986\ \Omega/\text{km}$.
  \item \textbf{Capacidad de corriente admisible (ambiente $40^\circ\text{C}$, viento $0.6\text{ m/s}$):} $145.0\text{ A}$.
  \item \textbf{Módulo de elasticidad final:} $6,000\text{ daN/mm}^2$.
  \item \textbf{Coeficiente de dilatación lineal:} $23.0 \times 10^{-6}\ ^\circ\text{C}^{-1}$.
\end{itemize}

\subsection{Cables Autoportantes de Baja Tensión: Tipo CAAI}
Para la red de distribución secundaria trifásica tetrafilar (380/220 V) y el alumbrado público se utilizarán cables concéntricos trenzados tipo CAAI (NTP 370.254 e IEC 60502-1), compuestos por fases aisladas con polietileno reticulado (XLPE resistente a intemperie y radiación UV) cableadas helicoidalmente alrededor de un portador neutro mensajero de aleación de aluminio desnudo o aislado:
\begin{table}[H]
\centering
\small
\begin{tabular}{|l|c|c|c|c|}
\hline
\textbf{Parámetro} & \textbf{Fases (3x35)} & \textbf{Neutro (NA 25)} & \textbf{AP (1x16)} & \textbf{Derivación (2x16)} \\
\hline
Conductor & Aluminio puro (EC) & Aleación Al (AAAC) & Aluminio puro (EC) & Aluminio puro (EC) \\
Sección nominal & $35.0\text{ mm}^2$ & $25.0\text{ mm}^2$ & $16.0\text{ mm}^2$ & $16.0\text{ mm}^2$ \\
Espesor aislamiento XLPE & $1.20\text{ mm}$ & Desnudo / $1.2\text{ mm}$ & $1.20\text{ mm}$ & $1.20\text{ mm}$ \\
Capacidad de corriente ($30^\circ\text{C}$) & $130.0\text{ A}$ & -- & $80.0\text{ A}$ & $80.0\text{ A}$ \\
Resistencia eléctrica ($20^\circ\text{C}$) & $0.868\ \Omega/\text{km}$ & $1.380\ \Omega/\text{km}$ & $1.910\ \Omega/\text{km}$ & $1.910\ \Omega/\text{km}$ \\
Carga rotura mensajero & -- & $7.40\text{ kN}$ & -- & $7.40\text{ kN}$ (NA 25) \\
\hline
\end{tabular}
\caption{Características técnicas de cables autoportantes CAAI en BT.}
\end{table}

\section{POSTES DE CONCRETO ARMADO CENTRIFUGADO (CAC)}

Los postes serán de concreto armado centrifugado de sección anular cónica, fabricados según la norma técnica NTP 339.027 y homologados por el MINEM:
\begin{itemize}
  \item \textbf{Postes de Media Tensión (10 kV):}
    \begin{itemize}
      \item CAC $12.0\text{ m} / 300\text{ daN}$: Para estructuras de alineamiento y ángulos pequeños ($< 15^\circ$). Diámetro de cima: 150 mm; base: 330 mm.
      \item CAC $13.0\text{ m} / 400\text{ daN}$: Para estructuras de anclaje, retención y subestaciones biposte. Cima: 180 mm; base: 375 mm.
      \item CAC $13.0\text{ m} / 600\text{ daN}$: Para terminales de línea en derivación de la MCH Limbani y cruce especial de quebrada.
    \end{itemize}
  \item \textbf{Postes de Baja Tensión (Red Secundaria y AP):}
    \begin{itemize}
      \item CAC $8.0\text{ m} / 200\text{ daN}$: Para alineamiento en pasajes peatonales y calles secundarias. Diámetro cima: 120 mm.
      \item CAC $8.0\text{ m} / 300\text{ daN}$: Para retención y cambio de dirección en calles secundarias.
      \item CAC $9.0\text{ m} / 300\text{ daN}$: Para avenidas principales, terminales de circuito y postes mixtos con soporte de luminarias de alumbrado público.
    \end{itemize}
\end{itemize}
Todos los postes contarán con orificios pasantes de 18 mm protegidos con tubos PVC para el paso de pernos de ferretería y orificio inferior para la puesta a tierra de armadura interior.

\section{TRANSFORMADORES DE DISTRIBUCIÓN}

Se especifican tres (03) transformadores trifásicos sumergidos en aceite mineral dieléctrico con refrigeración natural ONAN, aptos para operación a más de $3,050\text{ m.s.n.m.}$, cumpliendo IEC 60076 y NTP 370.002:
\begin{table}[H]
\centering
\small
\begin{tabular}{|l|c|c|c|}
\hline
\textbf{Característica} & \textbf{SED-01 (Norte)} & \textbf{SED-02 (Centro)} & \textbf{SED-03 (Sur)} \\
\hline
Potencia nominal & $50.0\text{ kVA}$ & $100.0\text{ kVA}$ & $100.0\text{ kVA}$ \\
Tensión primaria nominal & $10,000\text{ V}$ & $10,000\text{ V}$ & $10,000\text{ V}$ \\
Tensión secundaria nominal & $380 - 220\text{ V}$ & $380 - 220\text{ V}$ & $380 - 220\text{ V}$ \\
Grupo de conexión & Dyn5 (neutro accesible) & Dyn5 (neutro accesible) & Dyn5 (neutro accesible) \\
Frecuencia de trabajo & $60\text{ Hz}$ & $60\text{ Hz}$ & $60\text{ Hz}$ \\
Regulación de tensión (taps) & $\pm 2 \times 2.5\%$ & $\pm 2 \times 2.5\%$ & $\pm 2 \times 2.5\%$ \\
Tensión de cortocircuito ($U_{cc}$) & $4.0\%$ & $4.5\%$ & $4.5\%$ \\
Nivel de aislamiento (BIL primario) & $95\text{ kV}$ (derrateado $\ge 110\text{ kV}$) & $95\text{ kV}$ (derrateado $\ge 110\text{ kV}$) & $95\text{ kV}$ (derrateado $\ge 110\text{ kV}$) \\
Tipo de refrigeración & ONAN & ONAN & ONAN \\
Aceite dieléctrico & Mineral libre de PCB & Mineral libre de PCB & Mineral libre de PCB \\
Límite calentamiento aceite & $\le 50^\circ\text{C}$ sobre ambiente & $\le 50^\circ\text{C}$ sobre ambiente & $\le 50^\circ\text{C}$ sobre ambiente \\
\hline
\end{tabular}
\caption{Parámetros electromecánicos de los transformadores de distribución.}
\end{table}

\section{EQUIPOS DE MANIOBRA Y PROTECCIÓN EN MEDIA TENSIÓN}

\subsection{Seccionadores Cortacircuito Unipolares (Cut-Out)}
\begin{itemize}
  \item \textbf{Tensión nominal:} $15.0\text{ kV}$.
  \item \textbf{Nivel básico de aislamiento al impulso (BIL):} $110\text{ kV}$.
  \item \textbf{Corriente nominal continua:} $100\text{ A}$.
  \item \textbf{Capacidad de interrupción simétrica:} $8.0\text{ kA}$ eficaz.
  \item \textbf{Línea de fuga mínima:} $430\text{ mm}$ (apto para altitud $> 3,000\text{ m.s.n.m.}$).
  \item \textbf{Tubos portafusibles:} Fibra de vidrio con revestimiento interior extintor de arco.
  \item \textbf{Eslabones fusibles:} Tipo Chicote K (acción rápida), calibre normalizado de $3\text{ A}$ para SED-01 (50 kVA) y $6\text{ A}$ para SED-02 y SED-03 (100 kVA).
\end{itemize}

\subsection{Pararrayos de Óxido Metálico (ZnO)}
\begin{itemize}
  \item \textbf{Tecnología:} Resistencias no lineales de óxido de zinc (ZnO) sin explosores, poliméricos en goma de silicona de alta hidrofobicidad.
  \item \textbf{Tensión nominal ($U_r$):} $12.0\text{ kV}$.
  \item \textbf{Tensión de operación continua máxima (MCOV):} $10.2\text{ kV}$.
  \item \textbf{Corriente nominal de descarga ($8/20\ \mu\text{s}$):} $10.0\text{ kA}$ (Clase distribución pesada).
  \item \textbf{Línea de fuga específica:} $\ge 31\text{ mm/kV}$ de tensión máxima del sistema.
  \item \textbf{Desconectador automático:} Incorporado en la base con terminal de puesta a tierra flexible de cobre.
\end{itemize}

\section{AISLADORES Y FERRETERÍA DE LÍNEA}

\begin{itemize}
  \item \textbf{Aisladores de espiga (Pin):} Tipo ANSI 55-5 de porcelana vidriada marrón para estructuras de alineamiento en 10 kV.
  \item \textbf{Aisladores poliméricos tipo suspensión:} Núcleo de varilla de fibra de vidrio y faldones de caucho de silicona, tensión nominal 15 kV, carga mecánica mínima de falla $70\text{ kN}$, para estructuras de retención y anclaje.
  \item \textbf{Crucetas de madera tratada / perfiles de acero:} Crucetas de perfil de acero galvanizado en caliente de $75 \times 75 \times 6.4\text{ mm} \times 2.0\text{ m}$ o madera eucalipto tratada clase estructural.
  \item \textbf{Ferretería y accesorios metálicos:} Perno maquinado $5/8''$, perno de ojo, abrazaderas para poste de 180 mm y 330 mm, tuercas de ojo, grapas de anclaje tipo pistola de aleación de aluminio. Todo galvanizado en caliente según norma ASTM A153 / ASTM A123 con espesor de capa de zinc mínimo de $80\ \mu\text{m}$.
\end{itemize}

\section{TABLEROS DE DISTRIBUCIÓN Y ALUMBRADO PÚBLICO}

En cada una de las tres subestaciones se instalará un tablero de distribución adosado a uno de los postes de la SAB a una altura de $2.50\text{ m}$ sobre el nivel del terreno:
\begin{itemize}
  \item \textbf{Gabinete:} Fabricado en plancha de fierro galvanizado de $1.5\text{ mm}$ de espesor con pintura electrostática poliéster color gris RAL 7035, grado de protección intemperie IP55 e impacto IK10.
  \item \textbf{Interruptor principal de caja moldeada (MCCB):}
    \begin{itemize}
      \item SED-01 (50 kVA): MCCB tripolar regulable $3 \times 80\text{ A}$, poder de corte $I_{cu} \ge 25\text{ kA}$ a 380 V.
      \item SED-02 y SED-03 (100 kVA): MCCB tripolar regulable $3 \times 160\text{ A}$, poder de corte $I_{cu} \ge 25\text{ kA}$ a 380 V.
    \end{itemize}
  \item \textbf{Interruptores derivados de circuito:} Termomagnéticos bipolares y tripolares de $3 \times 50\text{ A}$, $3 \times 63\text{ A}$ y $2 \times 20\text{ A}$ en riel DIN.
  \item \textbf{Control de Alumbrado Público:} Interruptor termomagnético de $2 \times 20\text{ A}$, contactor electromagnético trifásico/bipolar de $32\text{ A}$ categoría AC-3 y fotocélula electrónica omnidireccional de encendido a 15 lux y apagado a 40 lux instalada en la cima del poste.
  \item \textbf{Totalizador de energía de subestación:} Medidor multifunción electrónico clase 0.5S con puerto de comunicación RS485 para monitoreo del balance de energía.
\end{itemize}

\section{LUMINARIAS DE ALUMBRADO PÚBLICO}

\begin{itemize}
  \item \textbf{Tecnología y fuente luminosa:} Luminaria LED integrada con módulos LED de alta eficiencia lumínica ($\ge 130\text{ lm/W}$).
  \item \textbf{Potencia consumida:} $50.0\text{ W}$ nominal. Flujo luminoso útil: $\ge 6,500\text{ lúmenes}$.
  \item \textbf{Temperatura de color correlacionada (CCT):} $4,000\text{ K}$ (blanco neutro para excelente visibilidad en neblina y lluvia).
  \item \textbf{Índice de reproducción cromática (CRI):} $\ge 70$.
  \item \textbf{Carcasa y disipador térmico:} Inyección de aleación de aluminio con aletas disipadoras y pintura anticorrosiva marina, grado IP66 en bloque óptico y compartimento de driver.
  \item \textbf{Driver electrónico:} Entrada de $120 - 277\text{ V}$, $60\text{ Hz}$, factor de potencia $\ge 0.95$, protección contra sobretensiones transitorias de $10\text{ kV} / 10\text{ kA}$.
  \item \textbf{Pastoral de soporte:} Tubo de acero galvanizado en caliente de $1 1/2'' \times 1.50\text{ m}$ con inclinación de $15^\circ$ respecto a la horizontal.
\end{itemize}

\section{MATERIALES PARA PUESTA A TIERRA}

\begin{itemize}
  \item \textbf{Varillas de puesta a tierra (jabalinas):} Acero revestido de cobre electrolítico depositado al $99.9\%$ (Copperweld) de $5/8''$ de diámetro por $2.40\text{ m}$ de longitud, con extremo cónico para clavado y cabeza recta para conector.
  \item \textbf{Conductor de enlace a tierra:} Alambre o cable concéntrico de cobre electrolítico temple blando desnudo de $35\text{ mm}^2$ (para neutro y masa de SED) y $25\text{ mm}^2$ (para finales de línea BT).
  \item \textbf{Conectores mecánicos:} Tipo grapa de bronce con perno de apriete tipo U o tipo AB de alta conductividad para conexión firme electrodo--cable.
  \item \textbf{Compuesto mejorador de resistividad:} Bentonita sódica activada granular de baja resistividad (presentación sacos de 25 kg) no corrosiva y respetuosa con el medio ambiente, complementada con sales minerales electrolíticas para reducir la resistividad del terreno rocoso.
  \item \textbf{Caja de registro:} De concreto prefabricado de $0.30 \times 0.30 \times 0.30\text{ m}$ con tapa de concreto y orificio central para inspección periódica con telurómetro.
\end{itemize}

\section{ACOMETIDAS DOMICILIARIAS Y SISTEMA DE MEDICIÓN}

\begin{itemize}
  \item \textbf{Cable concéntrico de acometida:} Cable bipolar concéntrico tipo SET de $2 \times 4.0\text{ mm}^2$ con conductores de cobre recocido y aislamiento de polietileno reticulado XLPE bajo cubierta de PVC resistente a rayos UV.
  \item \textbf{Caja portamedidor:} Cajas poliméricas monofásicas tipo L-T en policarbonato reforzado con fibra de vidrio autoextinguible (grado de protección IP54) con visor transparente de policarbonato anti-impacto.
  \item \textbf{Medidor de energía activa:} Medidor estático de energía activa monofásico clase 1.0 (IEC 62053-21), bifilar, tensión $220\text{ V}$, $60\text{ Hz}$, corriente nominal $5(60)\text{ A}$, display digital LCD con registro de energía acumulada en kWh y memoria no volátil.
  \item \textbf{Interruptor de protección en caja de acometida:} Termomagnético bipolar $2 \times 16\text{ A}$, curva C, poder de ruptura $I_{cn} \ge 6.0\text{ kA}$ a 220 V.
\end{itemize}
